\documentclass[10pt,a4paper]{amsart}
\usepackage[margin=19mm]{geometry}
\usepackage{amsmath,amssymb,mathtools}
\usepackage[scr=rsfs,cal=euler]{mathalfa}
\usepackage{xfrac}
\usepackage[unicode,hidelinks,bookmarks=false]{hyperref}
\newcommand{\ie}{i.e.\ }
\newcommand{\cf}{cf.\ }
\newcommand{\resp}{respectively\ }
\newcommand{\Subsection}{\subsection}
\providecommand{\nobreakdash}{\nobreak\hbox{-}}
\setlength{\emergencystretch}{2em}
\multlinegap0pt
\usepackage{amssymb}
\usepackage{dsfont}
\usepackage{mathtools}
\usepackage{xcolor}
\usepackage[shortlabels]{enumitem}
\newtheorem{lemma}{Lemma}
\newtheorem{theorem}{Theorem}
\DeclareMathOperator{\m}{M}
\title{$\m$-ideals: from Banach spaces to rings}
\author{David P. Blecher, Amartya Goswami}
\date{Source: arXiv:2504.14829v3 (CC BY 4.0)}
\begin{document}
\maketitle

\noindent\emph{Source and licence.} $\m$-ideals: from Banach spaces to rings. Version arXiv:2504.14829v3 (CC BY 4.0), distributed by arXiv under the Creative Commons Attribution 4.0 International licence, http://creativecommons.org/licenses/by/4.0/, as recorded in the arXiv licence metadata for this exact version and shown on the abstract page https://arxiv.org/abs/2504.14829v3. Full licence notice: \texttt{licenses/arxiv-2504.14829-CC-BY-4.0.txt}.\par\smallskip

\noindent\textit{Editorial scope.} Counted unit: the article's theorem fix (source0.tex lines 540--552). The article's complete original proof of it is delivered with it (source0.tex lines 554--586, one \texttt{proof} environment of 33 lines). No mathematics is written, repaired or re-derived: the statement and the proof are the article's own lines, in the article's own notation, and no retained line is substituted or re-flowed.  Retained uncounted as the source-local context the proof needs: lines 531--534.  Nothing was omitted from any retained span, so the package carries no omission record.  No retained line contains a citation, so the wrapper carries no bibliography and no bibliography entry.  No retained line contains a figure environment, an image-inclusion command or drawing code, so no image is delivered and the package declares no asset dependency.  Editorial formatting only, in addition: an AMS \texttt{amsart} preamble that restates the article's own declarations and macros used by the retained lines, namely the environments \texttt{lemma}, \texttt{theorem} on separate counters, together with the standard packages those lines need.  The retained environments print in this document as Lemma 1, Theorem 1 in order of appearance, whereas the article prints the same items with the numbering of its own full document.  The complete native source of the article as distributed in the arXiv e-print for this exact version is bundled unchanged as \texttt{sources/176829/source0.tex}.
\begin{lemma}
\label{pros}
Suppose that a unital ring $R$ has a nontrivial ideal decomposition $R = I \oplus J$.  Then $I$ is an $\m$-ideal in $R$ if and only if $I$ is relatively irreducible.
\end{lemma}

\begin{theorem} \label{fix} Let $n \geq 2$, and $R_1,  \cdots R_n$ be a finite collection of unital rings, with ideals $I_1, \cdots , I_n$ respectively, with
at least one of these is nonzero.
Let  $R =  \prod_{k=1}^n \, R_k$, the direct product, and  $I = \prod_{k=1}^n \, I_k$.
In places below we will  identify $I = \oplus_k \, I_k$, and think of this as a sum of ideals of $R$. 
\begin{enumerate}
\item If all the $I_k$ are nonzero then $I$ is  an $\m$-ideal in $R$ if and only if $I_k$ is essential in $R_k$ for all $k$.
\item Suppose that  $I_k = 0$ for some $k$.  Then   $I$ is an $\m$-ideal in $R$ if and only if $I = I_j$ for some $j$
and  $I_j$ is relatively irreducible.

\item Any  nonzero
$\m$-ideal of $R$ is such a direct product of $\m$-ideals, that is it is of the type in {\rm (1)} or {\rm 
 (2)} above. 
\end{enumerate} \end{theorem}

\begin{proof}  This relies on the well-known fact that the ideals of a direct product of a finite number of rings are exactly the direct products of
ideals.   
 
If $I$ is an $\m$-ideal and at least two of the $I_k$ are nonzero then $I$ is essential (from which it follows that all of the $I_k$ are nonzero and essential in $R_k$). 
For otherwise, if $J \cap I= 0$ for  ideal $J \neq 0$,
then $J + I_1$ and $J + I_2$ intersect each other nontrivially, and each intersects $I$ if $I_1$ and $I_2$ are
nonzero, but
their intersection does not  intersect $I$ nontrivially.  
 
(1)\ Suppose that $I_k$ is essential in $R_k$ for all $k$ and that
$J_k$ and $K_k$ are ideals  in $R_k$ for all $k$, with $J_i \cap K_i \neq 0$ say.  Then
$I_i \cap (J_i \cap K_i) \neq 0$.  So $I \cap (\prod_{i=1}^n \, J_i) \cap (\prod_{i=1}^n \, K_i) \neq 0$. Hence
$I$ is an $\m$-ideal.
 
Conversely, if one of the $I_k$ is not essential in $R_k$
then $I$ is not an $\m$-ideal.  For example, suppose that $I_1$ is not essential
with  $J \cap I_1 = 0$ for an ideal $J$ in $R_1$. Then $J \cap I = J \cap I_1 = 0$,
which contradicts the fact in the second paragraph of the proof if $I$ is an $\m$-ideal.

(2)\ If at least two of the $I_k$ are nonzero  then  they must be all nonzero to get an  $\m$-ideal
by the second paragraph.   If all the $I_k$ are zero except one, then 
the proof of  Lemma \ref{pros} gives (2). 

 
(3)\ Suppose that $I$  is a nonzero $\m$-ideal of $\prod_{i=1}^n \, 
R_i$.  Then by the fact in the first paragraph of the proof,
 $I =  \prod_{k=1}^n \, I_k$.
Suppose that $J_1$ and $K_1$ are ideals in $R_1$ with $J_1 \cap K_1, J_1 \cap I_1, I_1 \cap K_1$ nonzero.  Then the intersection of any two of
\[J_1\times 0\cdots\times 0,\; \; \; K_1 \times 0\cdots\times 0, \; \; \; I = I_1 \times \left(\prod_{i=2}^n \, I_k\right)\]
is nonzero.
Since $I$ is an $\m$-ideal, the intersection of these three is nonzero.  Thus $I_1 \cap J_1 \cap K_1 \neq 0$.
Hence $I_1$ is an $\m$-ideal, and similarly so is $I_2, \cdots, I_n$. Thus $I$ is a direct product of $\m$-ideals. 
\end{proof}

\end{document}
